\clearpage
\begin{summarybox}
\textbf{\textcolor{CSummary}{一句话抓住本质}}

\textbf{同样是连续对称两次，轴轴“不变号”，点点“变两次号”，所以都得到 $2d$；轴点“只变一次号”，要再来一遍，得到 $4d$。}

\medskip
\textbf{\textcolor{CSummary}{三条最短记忆}}
\[
\boxed{\begin{array}{c|c}
\text{两条轴} & T=2d\\[3pt]
\text{两个同高度的中心} & T=2d\\[3pt]
\text{一条轴和一个中心} & T=4d
\end{array}\qquad d=|a-b|>0.}
\]

\textbf{使用前检查：}（1）两个对称位置是否不同？（2）是轴还是点？（3）两个点的高度是否相同？（4）题目要的是\textbf{一个周期}还是\textbf{最小正周期}？

\medskip
\textbf{\textcolor{CSummary}{最值得记住的中间一步}}

轴点组合：
\[
\boxed{f(x+2(b-a))=2c-f(x)}
\quad\Longrightarrow\quad
\boxed{f(x+4(b-a))=f(x)}.
\]

\textbf{前半程相对高度反过来，后半程才回到原高度。}

\medskip
\textbf{\textcolor{CSummary}{30 秒自测}}
\begin{enumerate}[label=\arabic*.,leftmargin=*,itemsep=0.25em]
\item 轴在 $x=1$ 和 $x=4$：一个周期？\quad\textbf{答：$6$。}
\item 中心在 $(0,2)$ 和 $(5,2)$：一个周期？\quad\textbf{答：$10$。}
\item 轴在 $x=-1$，中心在 $(2,3)$：一个周期？\quad\textbf{答：$12$。}
\item 两中心在 $(0,0)$、$(2,2)$：能直接认定周期 $4$ 吗？\quad\textbf{答：不能。}
\end{enumerate}

\textbf{关联知识：}F020（对称轴判定）、F021（对称中心判定）、F022（周期等式）、F023--F024（有关周期与正负变化的关系）。
\end{summarybox}
